% ----------------------------------------------------------
\begin{frame}{Algebraic Laws --- Why They Matter}
  A \textbf{query optimiser} inside a DBMS never executes the query as written.
  Instead it \textbf{rewrites} the expression using algebraic equivalences to
  find a cheaper execution plan.

  \begin{block}{Core Principle}
    Two expressions are \textbf{equivalent} if they produce the same result
    for every valid database instance.  The DBMS may freely substitute one for
    the other.
  \end{block}

  \begin{examples}{Motivating Example}
    Consider
    $\rasel_{\texttt{Country}=\texttt{'DE'}}(\texttt{Customers} \times \texttt{Orders})$\\[4pt]
    The Cartesian product alone could produce \emph{millions} of rows.  The
    optimiser rewrites this by \textbf{pushing} the restriction
    \emph{inside} the product:
    \[
      \rasel_{\texttt{Country}=\texttt{'DE'}}(\texttt{Customers}) \times \texttt{Orders}
    \]
    Now only German customers participate in the product --- a dramatic
    reduction in intermediate-result size.
  \end{examples}
\end{frame}

% ----------------------------------------------------------
\begin{frame}{Laws: Commutativity}
  An operator is \textbf{commutative} if swapping the operands yields the
  same result.

  \begin{block}{Commutative Operators}
    \begin{itemize}
      \item \textbf{Union:}             $A \cup B = B \cup A$
      \item \textbf{Intersection:}      $A \cap B = B \cap A$
      \item \textbf{Cartesian Product:} $A \times B = B \times A$
      \item \textbf{Join:}              $A \rajoin B = B \rajoin A$
    \end{itemize}
  \end{block}

  \begin{alertblock}{Set Difference is NOT commutative}
    $A \setminus B \neq B \setminus A$ in general.
  \end{alertblock}

  \begin{block}{Practical Use}
    The DBMS may reorder join operands so that the \emph{smaller} table is
    processed first, reducing memory usage for hash joins.
  \end{block}
\end{frame}

% ----------------------------------------------------------
\begin{frame}{Laws: Associativity}
  An operator is \textbf{associative} if the grouping of operands does not
  matter.

  \begin{block}{Associative Operators}
    \begin{itemize}
      \item \textbf{Union:}
            $(A \cup B) \cup C = A \cup (B \cup C)$
      \item \textbf{Intersection:}
            $(A \cap B) \cap C = A \cap (B \cap C)$
      \item \textbf{Cartesian Product:}
            $(A \times B) \times C = A \times (B \times C)$
      \item \textbf{Join:}
            $(A \rajoin B) \rajoin C = A \rajoin (B \rajoin C)$
    \end{itemize}
  \end{block}

  \begin{block}{Practical Use --- Join Ordering}
    With commutativity and associativity together, a DBMS can try
    \textbf{all possible join orderings} for a multi-table query and pick
    the cheapest.  For $n$ tables there are $\frac{(2(n-1))!}{(n-1)!}$
    possible orderings --- a key part of cost-based optimisation.
  \end{block}
\end{frame}

% ----------------------------------------------------------
\begin{frame}{Laws: Restriction Rules \textbf{(!)}}
  \begin{alertblock}{Note: Not Required for Exam}
    The following equivalences are important for understanding query
    optimisation but will not be examined.
  \end{alertblock}

  Restrictions are \textbf{cheap} (they reduce tuple counts) and should be
  applied \emph{as early as possible} --- this is called \textbf{predicate
  push-down}.

  \begin{block}{Key Equivalences}
    \begin{itemize}
      \item \textbf{Merge:}
            $\rasel_{C_1}(\rasel_{C_2}(R)) = \rasel_{C_1 \wedge C_2}(R)$
      \item \textbf{Commutativity:}
            $\rasel_{C_1}(\rasel_{C_2}(R)) = \rasel_{C_2}(\rasel_{C_1}(R))$
      \item \textbf{Distribute over union:}
            $\rasel_C(R_1 \cup R_2) = \rasel_C(R_1) \cup \rasel_C(R_2)$
      \item \textbf{Distribute over join (push-down):}
            $\rasel_C(R_1 \rajoin R_2) = \rasel_C(R_1) \rajoin \rasel_C(R_2)$
      \item \textbf{Commute with projection:}
            $\rasel_C(\raproj_A(R)) = \raproj_A(\rasel_C(R))$
            \quad (restrict \emph{before} projecting)
    \end{itemize}
  \end{block}
\end{frame}

% ----------------------------------------------------------
\begin{frame}{Laws: Predicate Push-Down --- Worked Example \textbf{(!)}}
  \begin{alertblock}{Note: Not Required for Exam}
  \end{alertblock}

  \begin{block}{Scenario}
    \textbf{Goal:} Find the names of German customers and their order totals.

    \textbf{Naive plan:}
    \[
      \raproj_{\texttt{Name, Total}}\!\left(
        \rasel_{\texttt{Country}=\texttt{'DE'}}(\texttt{Customers} \rajoin \texttt{Orders})
      \right)
    \]
    The join produces a huge intermediate result before filtering.
  \end{block}

  \begin{block}{Optimised plan (push restriction inside join)}
    \[
      \raproj_{\texttt{Name, Total}}\!\left(
        \rasel_{\texttt{Country}=\texttt{'DE'}}(\texttt{Customers})
        \;\rajoin\;
        \texttt{Orders}
      \right)
    \]
    Filter \texttt{Customers} first $\Rightarrow$ only German customers
    participate in the join $\Rightarrow$ much smaller intermediate result.
  \end{block}

  Both expressions are \textbf{equivalent} --- the result is identical.
\end{frame}

% ----------------------------------------------------------
\begin{frame}{Laws: Division and Further Rules \textbf{(!)}}
  \begin{alertblock}{Note: Not Required for Exam}
    These rules deepen understanding of the algebra but are not examinable.
  \end{alertblock}

  \begin{block}{Push a Condition Into a Join Operand}
    If $C$ references only attributes of $R_2$:
    \[
      \rasel_{C}(R_1 \rajoin R_2) \;=\; R_1 \rajoin \rasel_{C}(R_2)
    \]
    Filtering $R_2$ \emph{before} the join reduces the number of tuples
    that participate.
  \end{block}

  \begin{block}{Division as the Inverse of Cartesian Product}
    \[
      R_1 = (R_1 \times R_2) \div R_2
    \]
    Just as $x = (x \cdot y) / y$ in arithmetic, division undoes a Cartesian
    product.  This identity confirms the intuitive meaning of the division
    operator.
  \end{block}

  \begin{block}{Projection Cascade}
    $\raproj_{A}(\raproj_{AB}(R)) = \raproj_{A}(R)$
    \quad (the outer projection subsumes the inner one)
  \end{block}
\end{frame}

% ----------------------------------------------------------
\begin{frame}{Laws: Summary for the Query Optimiser \textbf{(!)}}
  \begin{alertblock}{Note: Not Required for Exam}
  \end{alertblock}

  \begin{center}
    \small
    \begin{tabular}{lp{6.5cm}}
      \toprule
      \textbf{Law} & \textbf{Optimiser action} \\
      \midrule
      Commutativity of join    & Reorder tables so smaller table is on the left \\
      Associativity of join    & Reorder multi-table joins to minimise intermediate sizes \\
      Predicate push-down      & Move $\rasel$ as close to base tables as possible \\
      Projection push-down     & Move $\raproj$ inward to reduce column count early \\
      Merge restrictions       & Combine multiple $\rasel$ into one filter pass \\
      \bottomrule
    \end{tabular}
  \end{center}

  \vspace{0.6em}
  \begin{block}{The Big Picture}
    These equivalences are \textbf{not about correctness} --- the results are
    always the same.  They are about \textbf{efficiency}: reducing the number
    of rows and columns that must be materialised during query execution.
  \end{block}
\end{frame}
